% Test that a margin footnote is set once (didactic, noteoverflow=foot): its text
% runs once, as without didactic, whether it stays in the margin or goes to the
% foot, and no box is kept for it.  Built with \def\testcase{...}:
%   gls   glossaries' first use and a counter of TeX's stepped with \global in
%         two footnotes, nothing overflowing: the first note says "three-letter
%         acronym (TLA). BUMP1", the second "TLA. BUMP2";
%   many  3000 margin footnotes of one line, each beside a paragraph of three
%         lines, so that every note stays in the margin: the memory used is
%         close to the same document's with noteoverflow=carry (\testcarry),
%         and the log names no box allocated per note;
%   feet  1200 footnotes of seven lines or so in the margin, each beside a
%         paragraph of four lines: more than the margin holds, so that a few
%         hundred go to the foot (each says FOOTWORD, counted at the foot from
%         the PDF); the log names no box allocated per note and the run has no
%         error, and the memory used is in the log for the record.
% Build and checks: see the Makefile.
\ifdefined\testcase\else\def\testcase{gls}\fi
\documentclass{memoir}
\ifdefined\testcarry\PassOptionsToPackage{noteoverflow=carry}{didactic}\fi
\usepackage{didactic}
\def\glscase{gls}
\ifx\testcase\glscase
  \usepackage[acronym]{glossaries}
  \newacronym{tla}{TLA}{three-letter acronym}
\fi
\newcount\mycount
\newcount\notes
\def\feetcase{feet}
\newcommand\bump{\global\advance\mycount1 BUMP\the\mycount}
\begin{document}
\ifx\testcase\glscase
Text\footnote{First use: \gls{tla}. \bump} more text.

Text\footnote{Second use: \gls{tla}. \bump} more text.
\else\ifx\testcase\feetcase
\loop\ifnum\notes<1200
Some words of text\footnote{FOOTWORD A long note, of the length of a reference
in full, with authors, a title, the journal, its volume and pages and the year,
which takes several lines in the margin.} and more words of the paragraph to
fill it up, so that it takes four lines of the text, with the note beside it
taking more lines than that, so that the margin fills up and some of the notes
go to the foot of the page instead, as they do in a text that cites a lot.\par
\global\advance\notes 1 \repeat
\else
\loop\ifnum\notes<3000
Some words of text\footnote{A short note.} and more words of the paragraph to
fill it up, so that it takes three lines of the text, and the note beside it,
one line, has all the room it needs in the margin, as most notes do.\par
\global\advance\notes 1 \repeat
\fi\fi
\end{document}
